\section{Introduction}\label{sec:introduction}

Symplectic geometry begins with a contrast.  A symplectic transformation
preserves volume, but volume alone does not determine whether one domain can be
embedded into another by such a transformation.  Symplectic capacities assign
a scalar ``size'' to a domain while respecting symplectic embeddings, scaling,
and a common normalization.  They therefore turn part of this rigidity into a
quantitative question.

Viterbo proposed an isoperimetric principle for these quantities: among convex
domains of a fixed volume, the Euclidean ball should have maximal normalized
symplectic capacity \cite{Viterbo2000}.  For the Ekeland--Hofer--Zehnder
capacity in four dimensions, the form studied in this thesis is
\[
  \operatorname{sys}(K)
  =\frac{c_{\mathrm{EHZ}}(K)^2}{2\operatorname{vol}_4(K)}
  \leq 1.
\]
The numerator has the same scaling degree as the four-dimensional volume, so
the systolic ratio is unchanged by translations, positive scalings, and
linear symplectic transformations.  The conjectured inequality asserted that
no convex body could have a larger ratio than the ball.

The conjecture is false.  Haim--Kislev and Ostrover constructed a
four-dimensional counterexample from the Lagrangian product of two regular
pentagons in a suitable relative position \cite{HaimKislevOstrover2024}.  In
the conventions used here, $P_5$ is a centered regular pentagon and the
Lagrangian product places its two factors in the $q$- and $p$-coordinate
planes.  We use the symplectically equivalent representative
\[
  K_{\mathrm{HKO}}=P_5\times_L R(-\pi/2)P_5,
  \qquad
  \operatorname{sys}(K_{\mathrm{HKO}})
  =\frac{3+\sqrt5}{5}>1.
\]
This result ends the universal inequality, but not the geometric problem that
motivated it.  It leaves open how exceptional the counterexample is, whether
nearby or computationally generated polytopes improve it, and which features
of the minimizing dynamics account for its value.  Those questions are the
starting point of this thesis.

\subsection{Questions and scope}

The project probes the counterexample from several directions rather than
trying to force them into a single theorem.  Its central local question is:
is the HKO polytope already best among nearby ten-facet polytopes, once the
transformations that preserve the systolic ratio are removed?  A second
question is empirical: can systematic finite computation and standard
data-science methods find a new source of polytopes with systolic ratio larger
than $1$, or at least a reproducible rule for proposing such candidates?  A
third question concerns method: how can the relevant boundary dynamics be
reduced to finite calculations whose hypotheses and evidence boundaries
remain visible?  The proved profile of a structured one-parameter
pentagon family provides a smaller test of the same finite machinery.

The restrictions are deliberate.  Dimension four is the first symplectic
dimension in which the HKO example occurs and is the setting of the local
problem.  Polytopes replace a smooth boundary flow by finitely many facet
directions.  This nonsmoothness does not remove the dynamics: generalized Reeb
orbits still realize the EHZ capacity, and at least one minimizing orbit can be
chosen with a finite word of distinct facet directions.  Haim--Kislev's formula
then expresses the capacity as a global finite optimization over such words
\cite{HK2017}.  This finite formula was also one of the computational routes
used in the HKO work; the thesis continues that route while separating
numerical discovery from proof.  Low facet counts keep the search
accessible, while the ten-facet HKO body and Lagrangian products retain the
non-generic structures that motivate the investigation.  The thesis does not
claim that these restrictions capture all nearby convex bodies or all possible
search mechanisms.

\subsection{Principal results}

The central mathematical result is a local-maximality theorem for the HKO
body.  In the Hausdorff space of polytopes with exactly ten facets, every
sufficiently nearby body has systolic ratio at most
\((3+\sqrt5)/5\).  After the neighborhood is made smaller, equality occurs
precisely along the orbit generated by translations, positive scalings, and
linear symplectic transformations.  Thus the result is strict transverse to
the natural symmetries, but it is a theorem about the ten-facet stratum; it
does not assert local maximality when the facet count changes or among all
nearby convex bodies.

The proof combines a local chart with a finite algebraic certificate.  Nearby
ten-facet polytopes are represented by their labelled polar vertices, and a
$25$-dimensional slice removes the $15$ infinitesimal symmetry directions.
For $26$ selected facet words, smooth strictly positive feasible choices with
positive quadratic value give upper functions for the true systolic ratio that
all touch it at the HKO point.  Exact SageMath verification shows that their
derivative rows span the cotangent space of the slice and admit a strictly
positive convex relation summing to zero.  A finite upper-function argument
then gives a uniform first-order decrease in every nonzero transverse
\mbox{direction}.

\Needspace{3\baselineskip}
The Rust computation that found the finite choices is not a proof input: the
SageMath verifier reconstructs and checks the algebraic predicates used by the
argument.

Two finite descriptions of the capacity supply the common computational
foundation.  The first is the Haim--Kislev quadratic program, presented with a
careful distinction between a feasible word, a stationary candidate, a
fixed-word maximum, and the global maximum that computes the capacity.  The
second is an idealized exact real-arithmetic flow-graph search developed from
affine passages through facet-pair sections.  For normalized real facet data
satisfying explicit nondegeneracy conditions on transitions, short row lists,
and return maps, an exhaustive distinct-facet word search returns the EHZ
capacity.  A stronger collection of these conditions holds on an open dense
subset of every fixed real presentation chamber.  When the facet rows are
rational, the construction is an exact rational-arithmetic algorithm; the
current Rust implementation targets this rational-input corollary under an
additional caller contract.  The theorem uses the local flow-map geometry of
Chaidez--Hutchings \cite{CH2021}, while completeness follows from the
simple-minimizer theorem.  It is not an algorithm for every polytope, and it is
not identified with the full Chaidez--Hutchings Type~1/Type~2 analysis.

For local variation, the thesis derives the volume and fixed-word capacity
derivatives in normalized facet coordinates.  A nondegenerate constrained
maximum yields a smooth branch and its gradient.  At singular data, a weaker
but more robust construction is often enough: any smooth feasible choice of
strictly positive word weights with positive quadratic value gives a smooth
upper function for the systolic ratio.  This is the device used in the HKO
proof.  The thesis does not claim a complete first-order branch catalogue for
arbitrary non-generic polytopes.  A retained twelve-start numerical panel shows
finite-step improvement by the associated search model, while its unfinished
trajectories provide no endpoint or local-maximality evidence.

The computational search has a different evidential status.  A retained table
contains $14336$ random polytopes and random Lagrangian products, none with
systolic ratio larger than $1$.  In-table models recover substantial
structure in the already computed target values, but this is explanation of a
finite table rather than candidate generation.  In a separate pre-evaluation
experiment, frozen scalar rules screened $100000$ generated random-product
candidates; the selected and baseline union comprised $1675$ evaluated
rows, again with no value larger than $1$.  These results provide a bounded
negative benchmark for the specified samples, generator, and rules.  They do
not show that random search cannot succeed, exhaust the data-science
repertoire, or establish a general nonexistence statement.

\Needspace{12\baselineskip}
As a side result, the same finite capacity machinery determines a formula for
a one-parameter family.  If
\(K_\theta=P_5\times_L R(\theta)P_5\) and
\(d(\theta)\) is the distance from \(\theta\) to the nearest multiple of
\(\pi/5\), then
\[
  \operatorname{sys}(K_\theta)
  =\frac{5+2\sqrt5}{10}\sec^2 d(\theta).
\]
The maximum in this family is exactly the HKO value.  The capacity comparison
is exhaustive: for this Lagrangian product, it suffices to consider words with
two or three blocks of each factor type.  A SageMath classification over an
ordered number field
shows that no competing branch has smaller action, and continuity fills the
finitely many exceptional parameters.  This is the resulting profile for the
rotated-pentagon family, not a classification of general polygon products or
of all minimizing orbits.

\subsection{How the parts fit together}

The shared dependency is the passage from dynamics to finite data.  The
preliminaries fix the common capacity, action, facet, and product conventions.
The subsequent polytope analysis proves that the capacity admits a simple
minimizing generalized Reeb orbit and separates the orbit's facet word and
dwell times
from the base-point condition placing it on the boundary.  Only after this
reduction do the quadratic-program and flow-graph methods appear.  They answer
different finite questions: the former compares global action values, whereas
the latter follows actual breakpoints through affine sections.

The first-order calculation supplies local models and upper functions for the
HKO certificate and delimits what numerical ascent can establish.  The
data-science search instead asks a broader discovery question with bounded
empirical evidence.  The rotated-pentagon profile returns to exact computation
in a structured family and independently uses the finite product enumeration.

The final parts distinguish evidence types.  Three-dimensional stereographic
views are exploratory, not proof evidence.  The numerics and availability
discussions separate binary64 search from exact verification and identify
material retained for audit or continuation.  The later discussion of AI in
the research process is distinct from the factual disclosure and from the
mathematical contributions above.

\clearpage
